\section{随机变量的独立与依赖}
\label{sec:05-Probability/05-Joint-Distributions/02}

两只股票。你算了一下，\(P(X>0)=0.55\)，\(P(Y>0)=0.60\)。然后你又算了 \(P(X>0,\,Y>0)=0.33\)，发现 \(0.33 \approx 0.55 \times 0.60\)。你觉得不错，两只股票似乎互不影响。

然后你换了另一组阈值。\(P(X>0.02,\,Y>0.02)=0.08\)，而 \(P(X>0.02)P(Y>0.02)=0.12\)。不相等了。一两个事件的巧合不能说明任何问题。

事件独立要求的是：不管你看哪两个事件，联合概率都必须等于各自概率的乘积。随机变量的独立把这个要求推到了极致——不是某一对事件，而是所有可能的事件组合。

\subsection{独立性的定义：所有事件同时过关}

随机变量 \(X\) 与 \(Y\) 独立，定义为：对任意合适的集合 \(A,B\)，

\[
P(X\in A,\,Y\in B)
=P(X\in A)\,P(Y\in B).
\]

``任意集合''这三个字分量很重。它意味着你检查 \(X>0\) 和 \(Y>0\) 独立还不够；你还要检查 \(X>0.001\) 和 \(Y<-0.5\)、\(X\in[0.1,0.2]\) 和 \(Y\in[0.3,0.5]\)……无穷多对。这看起来是个不可能完成的任务。

好在，如果联合 PMF 或联合密度存在，独立性有一个可操作的等价条件：分解。

\[
p_{X,Y}(x,y)=p_X(x)\,p_Y(y)
\]

或

\[
f_{X,Y}(x,y)=f_X(x)\,f_Y(y)
\]

在相应意义下几乎处处成立。为什么这等价于对所有事件独立？因为任何事件区域上的概率都是密度（或 PMF）在该区域上的积分（或求和），联合密度分解后二重积分自然分离。所以检查联合 PMF 的每一个格子是否等于行和乘列和，或者检查联合密度是否在支撑上等于边缘密度的乘积，就足够了。

\subsection{支撑区域本身就能否定独立}

上一节的三角形支撑

\[
0<y<x<1
\]

提供了一个快捷判断：X 和 Y 各自的边缘支撑都是 \((0,1)\)——单独看每个变量，0 到 1 之间的任何取值都可以出现。但组合起来呢？比如 \(x=0.2,y=0.8\)：X 可以取 0.2，Y 可以取 0.8，但 \((0.2,0.8)\) 这个组合却不可能出现。三角形支撑在左上角缺了一大块。

这就直接判了死刑：如果 X 和 Y 独立，联合支撑必须是两个边缘支撑的笛卡尔积——一个矩形。三角形的联合支撑不是矩形，所以 X 和 Y 不可能独立。

这个几何判据很好用：看一眼联合支撑的形状。边界如果不是坐标轴对齐的矩形，大概率已经存在依赖。不需要计算任何条件分布或相关系数，从定义域就可以排除独立性。

\subsection{独立时条件分布退化回边缘分布}

如果 X 和 Y 独立，把联合密度分解代入条件密度的定义：

\[
f_{X\mid Y}(x\mid y)
=\frac{f_X(x)\,f_Y(y)}{f_Y(y)}
=f_X(x).
\]

知道了 Y 的值之后，X 的分布跟不知道 Y 时完全一样。这个结论和事件独立的直觉完美对接：如果两件事独立，知道一件事发生不改变另一件事的概率。

反过来想更深刻：如果你发现，不管 Y 取什么值，X 的条件分布总是同一个边缘分布，那联合分布就一定是边缘分布的乘积。条件恒等于边缘，等价于独立性。所以独立性不是抽象概念，是可以从条件分布中读出来的：你固定 Y 的不同取值，画几张 X 的条件分布图，如果每一张都一样，那 X 和 Y 就是独立的。

\subsection{独立让函数的期望也分解}

如果 \(X \perp Y\)（这个记号 \(\perp\) 表示相互独立），那么你用 X 构造的任何函数 \(g(X)\)，和用 Y 构造的任何函数 \(h(Y)\)，也是独立的。因为 \(g(X)\) 落在某个集合的事件等价于 X 落在原像集合的事件，后者只依赖 X；同理 \(h(Y)\) 只依赖 Y。

当期望存在时，这带来一个极其实用的分解。连续情形：

\[
\begin{aligned}
\E[g(X)h(Y)]
&=\iint g(x)h(y)\,f_X(x)f_Y(y)\,dx\,dy\\
&=\int g(x)f_X(x)\,dx \cdot \int h(y)f_Y(y)\,dy\\
&=\E[g(X)]\,\E[h(Y)].
\end{aligned}
\]

二重积分因为密度的因子分解而彻底分离成两个单重积分的乘积。取 \(g(X)=X,\;h(Y)=Y\)，就得到 \(\E[XY]=\E[X]\E[Y]\)，进而协方差 \(\Cov(X,Y)=0\)。但请留意方向：独立可以推出协方差为零，协方差为零推不回独立。独立要求的是对所有可测函数的分解，协方差只检查了二阶矩这一个函数对——这只是独立性的几千分之一。

\subsection{依赖不只是一条相关系数}

相关系数只捕捉了线性依赖。变量之间可以存在无限多种偏离独立的方式：

\begin{itemize}
\tightlist
\item
  线性正相关或负相关——相关系数能描述；
\item
  U 形依赖——Y 在 X 偏大和偏小时都高，在 X 中等时低。相关系数可能接近于零，但依赖非常明显；
\item
  尾部依赖——大部分时候两个变量各走各的，但极端事件（崩盘、峰值）同时出现。一阶二阶矩几乎察觉不到；
\item
  周期依赖——Y 随 X 周期性地涨落，但整体趋势不升不降；
\item
  隐变量混合——X 和 Y 本身独立，但都依赖第三个变量 Z。给定 Z 后 X 和 Y 条件依赖，但忽略 Z 看总体，X 和 Y 可能表现出虚假的关联或独立；
\item
  Simpson 悖论式反转——总体呈现正相关，但分组后每组都是负相关（或者反过来）。
\end{itemize}

联合分布完整记录了所有这些结构。相关系数只是联合分布在二阶矩上的一个投影，就像用手电筒照一个三维物体，墙上只留下一个影子。信息论中的互信息会提供更一般的依赖度量，能够捕捉非线性的关联，但它也只是一个数字——把整个依赖结构压缩成一个标量。理解依赖，最终还是要回到联合分布。

\subsection{从联合表格判断独立：一个具体计算}

设 X 和 Y 都只取 0 或 1，联合表格为：

\[
\begin{array}{c|cc}
&Y=0&Y=1\\
\hline
X=0&0.3&0.2\\
X=1&0.3&0.2
\end{array}
\]

先算行和列和。行和：\(X=0\) 为 \(0.3+0.2=0.5\)，\(X=1\) 也是 0.5。列和：\(Y=0\) 为 \(0.3+0.3=0.6\)，\(Y=1\) 为 \(0.2+0.2=0.4\)。

现在逐格检查分解：

\[
\begin{aligned}
P(X=0,Y=0)&=0.3,\qquad p_X(0)p_Y(0)=0.5\times0.6=0.3,\\
P(X=0,Y=1)&=0.2,\qquad p_X(0)p_Y(1)=0.5\times0.4=0.2,\\
P(X=1,Y=0)&=0.3,\qquad p_X(1)p_Y(0)=0.5\times0.6=0.3,\\
P(X=1,Y=1)&=0.2,\qquad p_X(1)p_Y(1)=0.5\times0.4=0.2.
\end{aligned}
\]

每一格都恰好等于行和乘列和。所以 \((X,Y)\) 独立。

现在想象你拿到的是另一张表。每一行和每一列跟上面完全一样（行 0.5/0.5，列 0.6/0.4），但内部四个数字不同——比如把右下角的 0.2 改成 0.3，那么为了总和仍为 1，至少另一个格子也要改。一旦任一格子不等于行和乘列和，独立性就破了。

这引出一个值得注意的点：不能只改动一个格子就讨论``依赖变强了还是变弱了''。改一个格子，边缘分布也可能跟着变，除非你同时调整其他格子保持行和列和不变。判断独立性必须看整张规范化表格，而不是孤立地盯一个数字。

下一节研究 \(S=X+Y\)。即使 X 和 Y 独立，S 的分布也不是两个密度逐点相乘，而是要汇总所有满足 \(x+y=s\) 的取值对——这引出一个新操作：卷积。
